\ifdefined\gkver\else\def\gkver{student}\fi
\documentclass[\gkver]{gaokaozhenti}
\begin{document}

\chapter{2004年北京理综}
%% paperType: 理综物理部分

\begin{enumerate}
\item
%% number: 14
%% paperName: 2004年北京理综第 14 题 6 分
%% typeId: 1
%% type: 单选题
%% chapter: 气体性质
%% point: 热力学第一定律
%% method: 无
%% score: 6
%% degree: 650
%% duplicateId: 0
%% body:
下列说法正确的是\xzanswer{D}
\fourchoices[answer=D]
{外界对气体做功，气体的内能一定增大}
{气体从外界吸收热量，气体的内能一定增大}
{气体的温度越低，气体分子无规则运动的平均动能越大}
{气体的温度越高，气体分子无规则运动的平均动能越大}
%% answer: D
%% memo:
\memoanswer{
本题原卷及材料中未提供详解，待补充。
}

\item
%% number: 15
%% paperName: 2004年北京理综第 15 题 6 分
%% typeId: 1
%% type: 单选题
%% chapter: 气体性质
%% point: 热力学第一定律
%% method: 无
%% score: 6
%% degree: 650
%% duplicateId: 0
%% body:
下列说法正确的是\xzanswer{D}
\fourchoices[answer=D]
{外界对气体做功，气体的内能一定增大}
{气体从外界吸收热量，气体的内能一定增大}
{气体的温度越低，气体分子无规则运动的平均动能越大}
{气体的温度越高，气体分子无规则运动的平均动能越大}
%% answer: D
%% memo:
\memoanswer{
由热力学第一定律 $\Delta U=W+Q$ 可知，$\Delta U$ 由 $W$ 和 $Q$ 共同决定，故 AB 错误；温度是分子平均动能的标志，故 C 错误 D 正确。
}

\item
%% number: 16
%% paperName: 2004年北京理综第 16 题 6 分
%% typeId: 1
%% type: 单选题
%% chapter: 机械波
%% point: 干涉衍射
%% method: 无
%% score: 6
%% degree: 650
%% duplicateId: 0
%% body:
声波属于机械波。下列有关声波的描述中正确的是\xzanswer{C}
\fourchoices[answer=C]
{同一列声波在各种介质中的波长是相同的}
{声波的频率越高，它在空气中传播的速度越快}
{声波可以绕过障碍物传播，即它可以发生衍射}
{人能辨别不同乐器同时发生的声音，证明声波不会发生干涉}
%% answer: C
%% memo:
\memoanswer{
声波的频率由波源决定，波速由介质决定，波长由波源和介质共同决定 $\left(\lambda=\frac{v}{f}\right)$，故 AB 错误；一切波都能发生干涉和衍射，故 C 正确 D 错误。
}

\item
%% number: 17
%% paperName: 2004年北京理综第 17 题 6 分
%% typeId: 1
%% type: 单选题
%% chapter: 原子和原子核
%% point: 玻尔模型
%% method: 无
%% score: 6
%% degree: 650
%% duplicateId: 0
%% body:
氦原子被电离一个核外电子，形成类氢结构的氦离子。已知基态的氦离子能量为 $E_{1}=-54.4\UeV$，氦离子能级的示意图如图所示。在具有下列能量的光子中，不能被基态氦离子吸收而发生跃迁的是\xzanswer{B}
\onepicture[w=4.0cm]{figs/17.png}
\fourchoices[answer=B]
{$40.8\UeV$}
{$43.2\UeV$}
{$51.0\UeV$}
{$54.4\UeV$}
%% answer: B
%% memo:
\memoanswer{
要吸收光子发生跃迁需满足一定条件，即吸收的光子能量必须是某两个能级的差值，由 $E_{2}-E_{1}=40.8\UeV$，$E_{4}-E_{1}=51.0\UeV$，$E_{\infty}-E_{1}=54.4\UeV$（电离），故 ACD 都可以，只有 B 不满足，故 B 正确。
}

\item
%% number: 18
%% paperName: 2004年北京理综第 18 题 6 分
%% typeId: 1
%% type: 单选题
%% chapter: 光学
%% point: 光的折射
%% method: 无
%% score: 6
%% degree: 650
%% duplicateId: 0
%% body:
已知一束可见光 a 是由 m、n、p 三种单色光组成的。检测发现三种单色光中，n、p 两种色光的频率都大于 m 色光；n 色光能使某金属发生光电效应，而 p 色光不能使该金属发生光电效应。那么，光束 a 通过三棱镜的情况是\xzanswer{A}
\fourchoices[answer=A, ispicture=true, h=2.2cm]
{figs/18a.png}
{figs/18b.png}
{figs/18c.png}
{figs/18d.png}
%% answer: A
%% memo:
\memoanswer{
由题意，n、p 色光的频率比 m 色光大，故 n、p 色光的折射率比 m 色光大，再由 n 色光能使某金属发生光电效应，而 p 色光不能，故 n 色光的折射率比 p 色光大，故 A 正确。
}

\item
%% number: 19
%% paperName: 2004年北京理综第 19 题 6 分
%% typeId: 1
%% type: 单选题
%% chapter: 磁场
%% point: 带电粒子在磁场中的运动
%% method: 无
%% score: 6
%% degree: 650
%% duplicateId: 0
%% body:
如图所示，正方形区域 abcd 中充满匀强磁场，磁场方向垂直纸面向里。一个氢核从 ad 边的中点 m 沿着既垂直于 ad 边又垂直于磁场的方向，以一定速度射入磁场，正好从 ab 边中点 n 射出磁场。沿将磁场的磁感应强度变为原来的 2 倍，其他条件不变，则这个氢核射出磁场的位置是\xzanswer{C}
\onepicture[w=4.6cm]{figs/19.png}
\fourchoices[answer=C]
{在 b、n 之间某点}
{在 n、a 之间某点}
{a 点}
{在 a、m 之间某点}
%% answer: C
%% memo:
\memoanswer{
带电粒子在磁场中的匀速圆周运动，由 $qvB=m\frac{v^{2}}{R}$ 得 $R=\frac{mv}{qB}$，且 a 点为圆心，若磁感应强度变为原来的 2 倍，半径变为原来的 $\frac{1}{2}$，所以氢核从 a 点射出，故 C 正确。
}

\item
%% number: 20
%% paperName: 2004年北京理综第 20 题 6 分
%% typeId: 1
%% type: 单选题
%% chapter: 曲线运动
%% point: 人造地球卫星
%% method: 无
%% score: 6
%% degree: 450
%% duplicateId: 0
%% body:
1990 年 5 月，紫金山天文台将他们发现的第 2752 号小行星命名为吴健雄星，该小行星的半径为 $16\Ukm$。若将此小行星和地球均看成质量分布均匀的球体，小行星密度与地球相同。已知地球半径 $R=6400\Ukm$，地球表面重力加速度为 $g$。这个小行星表面的重力加速度为\xzanswer{B}
\fourchoices[answer=B]
{$400g$}
{$\frac{1}{400}g$}
{$20g$}
{$\frac{1}{20}g$}
%% answer: B
%% memo:
\memoanswer{
设小行星和地球的密度均为 $\rho$，小行星的半径为 $r$、质量为 $m$，地球的半径为 $R$、质量为 $M$，星球表面重力加速度分别为 $g'$、$g$，物体质量为 $m_{0}$。在小行星表面有 $G\frac{mm_{0}}{r^{2}}=m_{0}g'$，$m=\rho\cdot\frac{4}{3}\pi r^{3}$；在地球表面有 $G\frac{Mm_{0}}{R^{2}}=m_{0}g$，$M=\rho\cdot\frac{4}{3}\pi R^{3}$。联立解得 $g'=\frac{1}{400}g$，故 B 正确。
}

\item
%% number: 21
%% paperName: 2004年北京理综第 21 题 6 分
%% typeId: 1
%% type: 单选题
%% chapter: 电场
%% point: 带电粒子的偏转
%% method: 无
%% score: 6
%% degree: 250
%% duplicateId: 0
%% body:
静电透镜是利用静电场使电子束会聚或发散的一种装置，其中某部分静电场的分布如下图所示。虚线表示这个静电场在 $xOy$ 平面内的一簇等势线，等势线形状相对于 $Ox$ 轴、$Oy$ 轴对称。等势线的电势沿 $x$ 轴正向增加，且相邻两等势线的电势差相等。一个电子经过 P 点（其横坐标为 $-x_{0}$）时，速度与 $Ox$ 轴平行。适当控制实验条件，使该电子通过电场区域时仅在 $Ox$ 轴上方运动。在通过电场区域过程中，该电子沿 $y$ 方向的分速度 $v_{y}$ 随位置坐标 $x$ 变化的示意图是\xzanswer{D}
\onepicture[w=6.0cm]{\includesvg{figs/21}}
\fourchoices[answer=D, ispicture=true, h=2.4cm]
{figs/21a.png}
{figs/21b.png}
{figs/21c.png}
{figs/21d.png}
%% answer: D
%% memo:
\memoanswer{
在 $x$ 轴负方向，电子所受电场力为右偏下，则电子的竖直速度沿 $y$ 轴负方向增加，到达 $O$ 点时最大，电子到达 $x$ 轴正方向时，所受电场力为右偏上，竖直分速度不断减小，由于电子一直向 $x$ 轴靠近，所经过区域在 $x$ 轴负向区电场线密，相应的加速度大，而在 $x$ 轴正向区加速度小，在经过相同的水平距离过程中，其竖直分速度不能减小到零，故 D 正确。
}

\item
%% number: 22
%% paperName: 2004年北京理综第 22 题 18 分
%% typeId: 4
%% type: 实验题
%% chapter: 电路
%% point: 电路实验
%% method: 无
%% score: 18
%% degree: 450
%% duplicateId: 0
%% body:
为了测定电流表 A$_{1}$ 的内阻，采用如图 \ref{2004北京理综22a} 所示的电路。其中：A$_{1}$ 是待测电流表，量程为 $300\UmuA$，内阻约为 $100\UO$；

A$_{2}$ 是标准电流表，量程是 $200\UmuA$；

$R_{1}$ 是电阻箱，阻值范围 $0\sim 999.9\UO$；

$R_{2}$ 是滑动变阻器；

$R_{3}$ 是保护电阻；

$E$ 是电池组，电动势为 $4\UV$，内阻不计；

S$_{1}$ 是单刀单掷开关，S$_{2}$ 是单刀双掷开关。

\onepicture[num=1, label=2004北京理综22a, w=5.0cm, align=right]{figs/22a.png}

\begin{enumerate}
	\item
	根据电路图 \ref{2004北京理综22a}，请在图 \ref{2004北京理综22b} 中画出连线，将器材连接成实验电路。

	\onepicture[num=2, label=2004北京理综22b, w=7.0cm]{figs/22b.png}

	\item
	连接好电路，将开关 S$_{2}$ 扳到接点 b 处，接通开关 S$_{1}$，调整滑动变阻器 $R_{2}$ 使电流表 A$_{2}$ 的读数是 $150\UmuA$；然后将开关 S$_{2}$ 扳到接点 a 处，保持 $R_{2}$ 不变，调节电阻箱 $R_{1}$，使 A$_{2}$ 的读数仍为 $150\UmuA$。若此时电阻箱各旋钮的位置如图 \ref{2004北京理综22c} 所示，电阻箱 $R_{1}$ 的阻值是 \tkanswer[1.5cm]{86.3} $\UO$，则待测电流表 A$_{1}$ 的内阻 $R_{g}=\tkanswer[1.5cm]{86.3}$ $\UO$。

	\onepicture[num=3, label=2004北京理综22c, w=3.6cm]{figs/22c.png}
	\item
	上述实验中，无论怎样调整滑动变阻器 $R_{2}$ 的滑动端位置，都要保证两块电流表的安全。在下面提供的四个电阻中，保护电阻 $R_{3}$ 应选用：\tkanswer[1.2cm]{B}（填写阻值相应的字母）。

	（A）$200\UkO$　　（B）$20\UkO$　　（C）$15\UkO$　　（D）$20\UO$
	\item
	下面提供最大阻值不同的四个滑动变阻器供选用。既要满足上述实验要求，又要调整方便，滑动变阻器 \tkanswer[1.2cm]{C}（填写阻值相应的字母）是最佳选择。

	（A）$1\UkO$　　（B）$5\UkO$　　（C）$10\UkO$　　（D）$25\UkO$
\end{enumerate}
%% answer:
\jdanswer{
\begin{enumerate}
	\item
	如解析图所示

	\item
	$86.3$，$86.3$

	\item
	B

	\item
	C
\end{enumerate}
}
%% memo:
\memoanswer{
（1）器材连接如图所示；

\onepicture[w=7.0cm]{figs/22_an.png}

（2）$86.3$；$86.3$；（3）B；（4）C。
}

\item
%% number: 23
%% paperName: 2004年北京理综第 23 题 18 分
%% typeId: 6
%% type: 计算题
%% chapter: 电磁感应
%% point: 电磁感应中的变加速
%% method: 无
%% score: 18
%% degree: 750
%% duplicateId: 0
%% body:
如图 \ref{2004北京理综23a} 所示，两根足够长的直金属导轨 MN、PQ 平行放置在倾角为 $\theta$ 的绝缘斜面上，两导轨间距为 $L$。M、P 两点间接有阻值为 $R$ 的电阻。一根质量为 $m$ 的均匀直金属杆 ab 放在两导轨上，并与导轨垂直。整套装置处于磁感应强度为 $B$ 的匀强磁场中，磁场方向垂直斜面向下。导轨和金属杆的电阻可忽略。让 ab 杆沿导轨由静止开始下滑，导轨和金属杆接触良好，不计它们之间的摩擦。

\onepicture[num=1, label=2004北京理综23a, w=6.5cm]{figs/23a.png}

\begin{enumerate}
	\item
	由 b 向 a 方向看到的装置如图 \ref{2004北京理综23b} 所示，请在此图中画出 ab 杆下滑过程中某时刻的受力示意图；

	\onepicture[num=2, label=2004北京理综23b, w=5.0cm]{figs/23b.png}
	\item
	在加速下滑过程中，当 ab 杆的速度大小为 $v$ 时，求此时 ab 杆中的电流及其加速度的大小；
	\item
	求在下滑过程中，ab 杆可以达到的速度最大值。
\end{enumerate}
%% answer:
\jdanswer{
\begin{enumerate}
	\item
	受力示意图如解析图所示（重力 $mg$，竖直向下；支持力 $N$，垂直斜面向上；安培力 $F$，沿斜面向上）

	\item
	$I=\frac{BLv}{R}$，$a=g\sin\theta-\frac{B^{2}L^{2}v}{mR}$

	\item
	$v_{m}=\frac{mgR\sin\theta}{B^{2}L^{2}}$
\end{enumerate}
}
%% memo:
\memoanswer{
（1）对 ab 杆受力分析如图所示，重力 $mg$，竖直向下；支持力 $N$，垂直斜面向上；安培力 $F$，沿斜面向上。

\onepicture[w=5.0cm]{figs/23_an.png}

（2）当 ab 杆速度为 $v$ 时，切割产生的感应电动势为 $E=BLv$，此时电路中电流为 $I=\frac{E}{R}=\frac{BLv}{R}$，ab 杆受到安培力为 $F=BIL=\frac{B^{2}L^{2}v}{R}$，由牛顿第二定律得 $mg\sin\theta-F=mg\sin\theta-\frac{B^{2}L^{2}v}{R}=ma$，解得 $a=g\sin\theta-\frac{B^{2}L^{2}v}{mR}$。

（3）当 ab 杆的向下的加速度为 0 时，速度达到最大速度 $v_{m}$，此时 ab 杆的合力为零，则有 $\frac{B^{2}L^{2}v}{R}=mg\sin\theta$，解得 $v_{m}=\frac{mgR\sin\theta}{B^{2}L^{2}}$。
}

\item
%% number: 24
%% paperName: 2004年北京理综第 24 题 20 分
%% typeId: 6
%% type: 计算题
%% chapter: 运动和力
%% point: 动量和能量
%% method: 无
%% score: 20
%% degree: 450
%% duplicateId: 0
%% body:
对于两物体碰撞前后速度在同一直线上，且无机械能损失的碰撞过程，可以简化为如下模型：A、B 两物体位于光滑水平面上，仅限于沿同一直线运动。当它们之间的距离大于等于某一定值 $d$ 时，相互作用力为零；当它们之间的距离小于 $d$ 时，存在大小恒为 $F$ 的斥力。设 A 物体质量 $m_{1}=1.0\Ukg$，开始时静止在直线上某点；B 物体质量 $m_{2}=3.0\Ukg$，以速度 $v_{0}$ 从远处沿该直线向 A 运动，如图所示。若 $d=0.10\Um$，$F=0.60\UN$，$v_{0}=0.20\Ums$，求：

\begin{enumerate}
	\item
	相互作用过程中 A、B 加速度的大小；
	\item
	从开始相互作用到 A、B 间的距离最小时，系统（物体组）动能的减少量；
	\item
	A、B 间的最小距离。
\end{enumerate}

\onepicture[w=7.0cm, align=right]{\includesvg{figs/24}}
%% answer:
\jdanswer{
\begin{enumerate}
	\item
	$a_{1}=0.60\Umsq$，$a_{2}=0.20\Umsq$

	\item
	$|\Delta E_{k}|=0.015\UJ$

	\item
	$\Delta s_{\min}=0.075\Um$
\end{enumerate}
}
%% memo:
\memoanswer{
（1）由题意知，在两物体相互作用的范围内，由牛顿第二定律得 $F=m_{1}a_{1}$，$F=m_{2}a_{2}$，代入解得 $a_{1}=0.60\Umsq$，$a_{2}=0.20\Umsq$。

\onepicture[w=6.0cm]{figs/24_an.jpg}

（2）两者速度相同时，相距最近，由动量守恒得 $m_{2}v_{0}=(m_{1}+m_{2})v$，解得 $v=\frac{m_{2}v_{0}}{(m_{1}+m_{2})}=0.15\Ums$，由能量守恒得 $|\Delta E_{k}|=\frac{1}{2}m_{2}v_{0}^{2}-\frac{1}{2}(m_{1}+m_{2})v^{2}=0.015\UJ$。

（3）由匀变速直线运动规律得 $v_{1}=a_{1}t$，$v_{2}=v_{0}-a_{2}t$，当 $v_{1}=v_{2}$ 时，解得 A、B 两者距离最近时所用时间 $t=0.25\Us$，由运动学公式得 $s_{1}=\frac{1}{2}a_{1}t^{2}$，$s_{2}=v_{0}t-\frac{1}{2}a_{2}t^{2}$，$\Delta s=s_{1}+d-s_{2}$，将 $t=0.25\Us$ 代入解得 A、B 间的最小距离 $\Delta s_{\min}=0.075\Um$。
}

\item
%% number: 25
%% paperName: 2004年北京理综第 25 题 22 分
%% typeId: 6
%% type: 计算题
%% chapter: 电场
%% point: 带电粒子的偏转
%% method: 无
%% score: 22
%% degree: 350
%% duplicateId: 0
%% body:
下图是某种静电分选器的原理示意图。两个竖直放置的平行金属板带有等量异号电荷，形成匀强电场。分选器漏斗的出口与两板上端处于同一高度，到两板距离相等。混合在一起的 a、b 两种颗粒从漏斗出口下落时，a 种颗粒带上正电，b 种颗粒带上负电。经分选电场后，a、b 两种颗粒分别落到水平传送带 A、B 上。

已知两板间距 $d=0.1\Um$，板的长度 $l=0.5\Um$，电场仅局限在平行板之间；各颗粒所带电量大小与其质量之比均为 $1\times10^{-5}\UCkg$。设颗粒进入电场时的初速度为零，分选过程中颗粒大小及颗粒间的相互作用力不计。要求两种颗粒离开电场区域时，不接触到极板但有最大偏转量。

\begin{enumerate}
	\item
	左右两板各带何种电荷？两极板间的电压多大？
	\item
	若两带电平行板的下端距传送带 A、B 的高度 $H=0.3\Um$，颗粒落至传送带时的速度大小是多少？
	\item
	设颗粒每次与传送带碰撞反弹时，沿竖直方向的速度大小为碰撞前竖直方向速度大小的一半。写出颗粒第 $n$ 次碰撞反弹高度的表达式。并求出经过多少次碰撞，颗粒反弹的高度小于 $0.01\Um$。
\end{enumerate}

\onepicture[w=5.0cm, align=right]{figs/25.png}
%% answer:
\jdanswer{
\begin{enumerate}
	\item
	$U=1\times10^{4}\UV$

	\item
	$v\approx 4\Ums$

	\item
	$h_{n}=\left(\frac{1}{4}\right)^{n}\times0.8\Um$；当 $n=4$ 时，$h_{n}<0.01\Um$
\end{enumerate}
}
%% memo:
\memoanswer{
（1）左板带负电荷，右板带正电荷。

依题意，颗粒在平行板间的竖直方向上满足
\[ l=\frac{1}{2}gt^{2}, \]①

在水平方向上满足
\[ s=\frac{d}{2}=\frac{1}{2}\cdot\frac{qU}{dm}t^{2}, \]②

①②两式联立得 $U=\frac{gmd^{2}}{2lq}=1\times10^{4}\UV$。

（2）颗粒落到水平传送带上，由动能定理得 $\frac{1}{2}qU+mg(l+H)=\frac{1}{2}mv^{2}$，解得 $v=\sqrt{\frac{qU}{m}+2g(l+H)}\approx 4\Ums$。

（3）在竖直方向颗粒做自由落体运动，它第一次落到水平传送带上沿竖直方向的速度为 $v_{y}=\sqrt{2g(l+H)}=4\Ums$，反弹的高度为 $h_{1}=\frac{(0.5v_{y})^{2}}{2g}=\left(\frac{1}{4}\right)\left(\frac{v_{y}^{2}}{2g}\right)$，由题设条件，颗粒第 $n$ 次反弹后上升高度为 $h_{n}=\left(\frac{1}{4}\right)^{n}\left(\frac{v_{y}^{2}}{2g}\right)=\left(\frac{1}{4}\right)^{n}\times0.8\Um$，当 $n=4$ 时，$h_{n}<0.01\Um$。
}

\end{enumerate}

\end{document}
